\documentclass[11pt, oneside]{article}
\usepackage[margin=1in]{geometry}
\geometry{letterpaper}
\usepackage{booktabs}
\usepackage{hyperref}
\usepackage{enumitem}
\usepackage{float}
\usepackage{caption}
\usepackage{subcaption}
\usepackage[numbers]{natbib}
\usepackage{graphicx}
\usepackage{amssymb}
\usepackage{amsmath}
\usepackage{array}
\usepackage{longtable}
\usepackage{xcolor}
\usepackage{tabularx}
\usepackage{setspace}
\hypersetup{
    colorlinks=true,
    linkcolor=blue,
    citecolor=blue,
    urlcolor=blue
}

\title{\textbf{CS781 Capstone Project:\\In-House Legal Intelligence System for Indian Jurisprudence}}
\author{
    \textbf{Abhinav S} \\
    Roll Number: \textbf{230041}
}
\date{}

\begin{document}
\maketitle

\begin{abstract}
Confidential, self-hosted Natural Language Processing (NLP) is essential for legal institutions handling sensitive judicial dockets subject to strict data-privacy and sovereignty mandates. This report details the architectural design, implementation, and systematic evolution of an end-to-end legal intelligence platform developed for Indian jurisprudence across seven distinct tasks: Foundation Pre-training, Legal NER, Rhetorical Role Prediction, Court Judgment Prediction, Legal Statute Identification, Abstractive Summarization, and Prior Case Retrieval. Constrained to a 50M--300M parameter ceiling with all external pre-trained checkpoints strictly prohibited, we architected and pre-trained from scratch a 210M-parameter 6-layer Hybrid Mixture-of-Experts (MoE) model on Indian court judgments, statutory Acts, and tribunal records. We transferred this foundation model across downstream tasks using parameter-efficient fine-tuning (PEFT) and knowledge distillation. Our findings demonstrate that sparse expert routing, task-specific document windowing, multiplicative parameter adaptation (LoRMA), and calibration overcome severe document length limitations and class imbalances, matching or approaching academic benchmarks under resource constraints.
\end{abstract}

% ============================================================
\section{Competition Result}
% ============================================================
\noindent
\textbf{Codabench username:} \underline{\hspace{3.5cm}}
\hfill
\textbf{Roll number:} \textbf{230041}
\medskip

\noindent
\textbf{Public GitHub Repository:} \url{https://github.com/[Your-Username]/CS781-Legal-NLP-Foundation}
\medskip

\noindent
Table~\ref{tab:leaderboard} summarizes the primary evaluation metrics and performance achieved by our best-performing models across all seven phases on the official test evaluation benchmark.

\begin{table}[H]
\centering
\small
\caption{Official evaluation results and primary performance metrics across all seven phases.}
\label{tab:leaderboard}
\begin{tabularx}{\textwidth}{clXc}
\toprule
\textbf{Phase} & \textbf{Task} & \textbf{Primary Evaluation Metric (Best Achieved)} & \textbf{Valid Runs} \\
\midrule
1 & Legal LM Pretraining & BERTScore F1: 0.7042 (ROUGE-L: 0.281, BLEU-4: 0.124) & 4 \\
2 & Legal NER & Strict Macro-F1: 0.3842 (Macro-P: 0.412, Macro-R: 0.360) & 6 \\
3 & Rhetorical Role Prediction & Strict Macro-F1: 0.3053 (Macro-P: 0.342, Macro-R: 0.281) & 6 \\
4 & Court Judgment Prediction & Strict Macro-F1: 0.5841 (F1\_0: 0.621, F1\_1: 0.547) & 6 \\
5 & Legal Statute Identification & Macro-F1: 0.2014 (Micro-F1: 0.312, Subset Acc: 0.145) & 6 \\
6 & Abstractive Summarization & BERTScore F1: 0.6814 (ROUGE-1: 0.321, ROUGE-L: 0.214) & 6 \\
7 & Prior Case Retrieval (PCR) & Micro-F1@7: 0.3500 (Micro-P@7: 0.344, Micro-R@7: 0.439) & 6 \\
\bottomrule
\end{tabularx}
\end{table}

\noindent
Across every phase, we satisfied the course requirement of testing at least 4--6 distinct architectural variations (including BitFit, Series Adapters, AdaLoRA, LoRMA, and Knowledge Distillation), ensuring systematic empirical comparisons rather than hyperparameter-only changes.

\section{Introduction and Problem Description}
% ============================================================
\noindent
Deploying language models within legal and judicial workflows presents unique challenges that off-the-shelf general-domain models cannot resolve. Indian court judgments are characterized by complex statutory cross-references, extensive document lengths (regularly exceeding 20--50 pages), archaic phraseology, extreme category imbalance, and strict confidentiality requirements. Sending privileged case records to third-party commercial APIs violates attorney-client privilege and judicial data sovereignty. Therefore, the core objective of this project was to engineer a sovereign, in-house legal intelligence platform trained strictly on official open legal corpora without relying on proprietary weights or external APIs.

The system was developed across seven interconnected judicial language tasks based on the Indian Legal Text Understanding and Reasoning (IL-TUR) benchmark~\cite{joshi2024iltur}: foundation language model pretraining from scratch (Phase~1), legal named entity recognition across 12 categories (Phase~2), rhetorical role discourse parsing across 13 roles (Phase~3), court judgment appeal outcome prediction (Phase~4), multi-label legal statute identification across 120 Acts (Phase~5), abstractive Supreme Court summarization (Phase~6), and prior case precedent retrieval across a 991-judgment repository (Phase~7).

Our design addressed three fundamental engineering challenges: (1) scaling pre-trained model capacity to 210M parameters while ensuring efficient single-GPU training; (2) evaluating parameter-efficient adaptation methods (LoRMA, AdaLoRA, BitFit) against Knowledge Distillation to retain pre-trained domain representations without catastrophic forgetting; and (3) preserving decisive judicial facts and operative holdings when documents vastly exceed transformer context windows.

\section{Data Analysis and Preprocessing}
% ============================================================
\noindent
All experimental data was sourced strictly from the official repository \texttt{Exploration-Lab/CS-781-Capstone}. Table~\ref{tab:corpus_stats} summarizes the corpus distributions, split sizes, and critical structural properties across all seven phases that informed our preprocessing pipelines and model architectures.

\begin{table}[H]
\centering
\small
\caption{Official corpus distribution, split sizes, and critical structural properties across phases.}
\label{tab:corpus_stats}
\begin{tabularx}{\textwidth}{clcccX}
\toprule
\textbf{Phase} & \textbf{Corpus Split / Task} & \textbf{Train Size} & \textbf{Dev Size} & \textbf{Test Size} & \textbf{Critical Structural Property} \\
\midrule
1 & Legal Pretraining Corpus & 4.2M tokens & 450K tokens & 237 docs & 4 Sources (SC, HC, Acts, Tribunals) \\
2 & Legal NER & 1,000 docs & 150 docs & 150 docs & 12 entity types; 86.4\% non-entity tokens \\
3 & Rhetorical Roles (CL/IT) & 24,180 sents & 3,240 sents & 3,110 sents & Extreme class imbalance; rare dissent roles \\
4 & Court Judgment (CJPE) & 37,387 docs & 4,154 docs & 4,155 docs & Class 0 (59.0\%) vs. Class 1 (41.0\%); long docs \\
5 & Legal Statute Ident. (LSI) & 4,821 docs & 602 docs & 603 docs & Multi-label; 120 statutes; long-tail power law \\
6 & Abstractive Summarization & 6,420 pairs & 800 pairs & 802 docs & Source avg. 5,120 words $\to$ summary 310 words \\
7 & Prior Case Retrieval (PCR) & 1,480 queries & 200 queries & 237 queries & Candidate pool of 991 cases; avg. 5.2 priors/query \\
\bottomrule
\end{tabularx}
\end{table}

\section{Phase-Wise Approach, Architectural Evolution, and Results}
% ============================================================
\noindent
For each phase, we outline the model architecture built, the rationale for selecting our final method over evaluated alternatives, the iterative debugging and failure analyses encountered, and the final top performance achieved.

% ------------------------------------------------------------
\subsection{Phase 1: Foundation Legal Language Model Pretraining}
% ------------------------------------------------------------
\noindent
\textbf{Methodology \& Architectural Selection:} We pre-trained an in-house domain foundation model from scratch within the required 50M--300M parameter limit, comparing a standard 12-layer Dense Transformer (124M parameters) against a custom 6-layer \textbf{Hybrid Mixture-of-Experts (HybridMoE)} architecture (210.45M total parameters; $d=768$, 12 heads, $T=512$, Byte-level BPE). Scaling a dense model to 210M parameters on a single GPU triggers out-of-memory errors due to quadratic matrix multiplication FLOPs ($O(L \cdot d_{\text{model}}^2)$). HybridMoE resolved this by alternating dense layers ($0, 2, 4$) with Top-1 MoE blocks ($1, 3, 5$) comprising 4 feed-forward experts. Because only one expert activates per token ($k=1$), per-token active FLOPs remain equivalent to a lightweight 124M model (124.03M active parameters) while total representational capacity expands to 210.45M parameters.

\medskip
\noindent
\textbf{Iterations, Failure Analysis \& Result:} Initial MoE runs suffered from expert collapse under Top-2 gating, where expert 1 handled $>80\%$ of tokens and starved remaining experts. We resolved this by implementing Top-1 gating with Gaussian exploration noise ($\epsilon \sim \mathcal{N}(0, \sigma^2)$) and an auxiliary load-balancing loss, achieving uniform expert utilization. This architectural enhancement dropped validation perplexity from 24.1 (dense baseline) down to 18.4, achieving a top score of \textbf{BERTScore F1: 0.7042} (ROUGE-L: 0.281, BLEU-4: 0.124).

% ------------------------------------------------------------
\subsection{Phase 2: Legal Named Entity Recognition (L-NER)}
% ------------------------------------------------------------
\noindent
\textbf{Methodology \& Architectural Selection:} To extract 12 legal entity classes mapped to 25 BIO tags under strict character-level evaluation, we evaluated Full Fine-Tuning (FFT), Linear Probing, BitFit, Additive LoRA ($r=16, \alpha=32$), \textbf{Low-Rank Multiplicative Adaptation (LoRMA)}, and a 3-layer compact student ($d=384, \sim20\text{M params}$). Under strict span matching, even a single-character boundary shift yields zero reward. We observed that additive LoRA ($W_0 + \frac{\alpha}{r}BA$) introduces additive perturbations that alter the geometric distances of pre-trained subword embeddings, inducing boundary flutter. In contrast, LoRMA modulates activations multiplicatively:
\begin{equation}
h = W_0 x + \alpha \cdot B(A(W_0 x))
\end{equation}
acting as an adaptive gain control that rescales representations along low-rank subspaces while preserving pre-trained token geometry. LoRMA outperformed all alternatives, reaching 0.3842 Strict Macro-F1 compared to 0.3421 for FFT and 0.3015 for BitFit.

\medskip
\noindent
\textbf{Iterations, Failure Analysis \& Result:} In early runs, standard Cross-Entropy loss yielded zero recall across minority classes (\texttt{WIT}, \texttt{R.COUNSEL}) because non-entity background tokens (\texttt{"O"}) comprise $86.4\%$ of the corpus. Replacing Cross-Entropy with an inverse-frequency power-weighted loss schedule ($w_c \propto (1/f_c)^{0.5}$) and zero-drift character offset refinement boosted minority detection by $+6.8\%$, securing our top score of \textbf{Strict Macro-F1: 0.3842} (Macro-P: 0.412, Macro-R: 0.360).

% ------------------------------------------------------------
\subsection{Phase 3: Rhetorical Role (RR) Prediction}
% ------------------------------------------------------------
\noindent
\textbf{Methodology \& Architectural Selection:} For classifying sentences into 13 judicial discourse roles across Corporate Law and Information Technology domains, we evaluated Linear Probing, BitFit, LoRA, Rank-Stabilized LoRA (rsLoRA), Student Distillation, and \textbf{LoRMA with 6D Document Flow Embeddings}. Because judicial discourse follows a canonical progression (Facts in opening recitals, Arguments in the body, and Rulings at closing), isolated sentence encoders struggle to distinguish between lower court recitals and present court rulings. We augmented the sentence representations with a 6D sinusoidal document depth vector $[\tau, \sqrt{\tau}, \sin(\pi \tau), \cos(\pi \tau), \sin(2\pi \tau), \cos(2\pi \tau)]$, where $\tau = i/(N-1)$ denotes the sentence position. This structural prior provided the model with global narrative trajectory.

\medskip
\noindent
\textbf{Iterations, Failure Analysis \& Result:} In our initial pipeline, inputs were formatted as \texttt{[PREV] \{prev\} [SENT] \{curr\}}. When the preceding context was lengthy, right-side truncation discarded the current sentence entirely, severely hurting classification. Inverting the sequence to target-first formatting (\texttt{\{curr\} [CONTEXT] \{prev\} \{next\}}) coupled with post-hoc logit adjustment ($\tilde{\mathbf{z}} = \mathbf{z} - \beta \log(\boldsymbol{\pi})$) restored recall on rare dissent roles, delivering our highest score of \textbf{Strict Macro-F1: 0.3053} (Macro-P: 0.342, Macro-R: 0.281).

% ------------------------------------------------------------
\subsection{Phase 4: Court Judgment Prediction and Explanation (CJPE)}
% ------------------------------------------------------------
\noindent
\textbf{Methodology \& Architectural Selection:} Predicting binary appeal outcomes ($0 = \text{Dismissed}, 1 = \text{Allowed}$) across 37,387 judgments requires processing texts averaging over 4,000 words. Naive truncation to the first 512 tokens discards judicial holdings entirely, plateauing at 0.51 Macro-F1, while full sliding-window chunking across 10,000 tokens is computationally prohibitive. We engineered a \textbf{Dual-Anchor Salient Windowing} module that extracts Anchor 0 (Head 512 tokens: charges and lower court history) and Anchor 1 (Tail 512 tokens: operative decree and ratio decidendi). Projecting both anchors into an interactive representation $[a_0; a_1; |a_1 - a_0|; a_1 \odot a_0] \in \mathbb{R}^{3072}$ enabled LoRMA to capture $95\%$ of verdict signal in under 8 minutes per training run.

\medskip
\noindent
\textbf{Iterations, Failure Analysis \& Result:} Initial training on the full 37,387 corpus induced majority-class bias toward dismissal (59\% class 0) and training instability. Introducing stratified subsampling (6,000 class 0 + 6,000 class 1) eliminated verdict bias and cut training steps from 9,347 to 750 batches. Combined with an empirical validation decision threshold sweep ($\theta^* = 0.48$), this improved Macro-F1 from 0.52 to our top performance of \textbf{Strict Macro-F1: 0.5841} ($F1_0: 0.621, F1_1: 0.547$).

% ------------------------------------------------------------
\subsection{Phase 5: Legal Statute Identification (LSI)}
% ------------------------------------------------------------
\noindent
\textbf{Methodology \& Architectural Selection:} To ground cases against 120 target Indian Acts, we compared Layer Freezing, BitFit, Bottleneck Adapters, AdaLoRA, and \textbf{LoRMA with Dynamic Per-Statute Thresholding}. Indian statutes follow an extreme power-law distribution: the top 5 provisions (e.g., IPC Section 302, CrPC Section 482) appear in $>20\%$ of cases, while over 80 statutes occur in $<0.5\%$ of cases. Standard binary classification with a fixed 0.5 threshold predicted only dominant statutory codes, yielding zero recall across rare acts. We solved this by implementing independent threshold calibration on validation data ($\theta_i^* \in [0.15, 0.65]$ for each statute $i$), enabling the model to surface rare, specific provisions without false positive inflation.

\medskip
\noindent
\textbf{Iterations, Failure Analysis \& Result:} Early experiments utilizing multi-class Softmax classification failed because court judgments routinely cite multiple statutes concurrently (averaging 3--7 provisions per petition). Transitioning to independent binary cross-entropy sigmoids over hierarchical chunk-pooled representations resolved multi-label co-occurrence dependencies, elevating performance to our top score of \textbf{Macro-F1: 0.2014} (Micro-F1: 0.312, Subset Accuracy: 0.145).

% ------------------------------------------------------------
\subsection{Phase 6: Abstractive Legal Summarization}
% ------------------------------------------------------------
\noindent
\textbf{Methodology \& Architectural Selection:} We formulated legal case summarization as conditional causal language modeling, comparing BitFit, Layer Freezing, Bottleneck Adapters, AdaLoRA, LoRMA, and a \textbf{Multi-Model Consensus Ensemble}. To compress lengthy judgments (averaging 5,120 words) into 300-word summaries within a 512-token budget, our architecture combined three mechanisms: (1) Salient Windowing (Lead 200 words + Operative Holding 120 words); (2) Conditional Loss Masking, which zeroed gradients over the input case prompt and optimized strictly over ground-truth summary tokens; and (3) Dynamic Frequency Penalties ($1.20 + 0.04 \times \text{count}$) to suppress repetitive judicial boilerplate.

\medskip
\noindent
\textbf{Iterations, Failure Analysis \& Result:} In early iterations, unconstrained greedy decoding generated pathological loops (e.g., repeatedly generating ``the appellant...''), while high temperatures induced factual hallucinations. We resolved this by pairing Nucleus sampling (top-$p=0.90$, temperature $0.70$) with an extractive fallback guard that triggered if generation produced fewer than 15 tokens. An ensemble combining LoRMA with length-normalized decoding achieved our top score of \textbf{BERTScore F1: 0.6814} (ROUGE-1: 0.321, ROUGE-L: 0.214).

% ------------------------------------------------------------
\subsection{Phase 7: Prior Case Retrieval (PCR)}
% ------------------------------------------------------------
\noindent
\textbf{Methodology \& Architectural Selection:} Given a query judgment, the objective is to retrieve and rank the top 7 relevant cited precedents from a 991-case repository. We evaluated a Dense Bi-Encoder, BitFit, Bottleneck Adapters, AdaLoRA, \textbf{Knowledge Distillation into a 42M Student}, and Multi-Model Reciprocal Rank Fusion (RRF). In precedent matching, our compact 42M student model achieved superior generalization over full-size teachers because distillation regularized the embeddings against overfitting to spurious lexical co-occurrences in lengthy judgments. Fusing student dense similarity with PEFT variants via weighted RRF ($w_{\text{distill}} = 2.5, k=60$) produced our optimal ranking system.

\medskip
\noindent
\textbf{Iterations, Failure Analysis \& Result:} Initial retrieval runs plateaued at 0.33 Micro-F1@7. Detailed error inspection uncovered a critical data condition: \textbf{in 100\% of test queries (237/237), the candidate pool contained the query case itself}. Because a case naturally has a cosine similarity of $1.0$ with itself, the model ranked the query case at Slot 1 across every query. Since a judgment cannot cite itself as a prior precedent, Slot 1 was a guaranteed false positive, and the true 7th precedent was truncated. Implementing an explicit self-exclusion filter ($c \neq q$) along with per-query Min-Max score normalization resolved the flaw and boosted performance to our top score of \textbf{Micro-F1@7: 0.3500} (Micro-Precision@7: 0.344, Micro-Recall@7: 0.439). In the published IL-TUR ACL 2024 benchmark~\cite{joshi2024iltur}, standard pre-trained baselines achieved single digits (InLegalBERT: 7.57\%, LegalBERT: 8.67\%); our system achieves a \textbf{$4.6\times$ improvement over InLegalBERT} and operates within 4\% of the world academic record (39.15\%).

\section{Experiments and Implementation Details}
% ============================================================
\noindent
All experiments were executed on single NVIDIA Tesla P100 (16GB) or T4 GPUs via Kaggle and Google Colab environments. Training utilized AdamW with decoupled weight decay ($\lambda=0.01$), linear warmup across $10\%$ of steps, cosine decay, and mixed-precision (FP16 via \texttt{torch.cuda.amp}).

\begin{table}[H]
\centering
\small
\caption{Hyperparameter specifications across all seven phases.}
\label{tab:hparams}
\begin{tabularx}{\textwidth}{lXXXXXXX}
\toprule
\textbf{Hyperparameter} & \textbf{P1} & \textbf{P2} & \textbf{P3} & \textbf{P4} & \textbf{P5} & \textbf{P6} & \textbf{P7} \\
\midrule
Total Parameters & 210.4M & 210.5M & 210.5M & 212.8M & 210.6M & 210.4M & 210.5M \\
Trainable Parameters & 210.4M & 2.5M & 2.5M & 1.8M & 2.5M & 2.5M & 2.5M \\
Max Context ($T$) & 512 & 512 & 160 & $2 \times 512$ & $4 \times 512$ & 512 & 512 \\
Peak Learning Rate & $3 \times 10^{-4}$ & $2 \times 10^{-4}$ & $3 \times 10^{-4}$ & $2 \times 10^{-4}$ & $3 \times 10^{-4}$ & $2 \times 10^{-4}$ & $3 \times 10^{-4}$ \\
Effective Batch Size & 32 & 16 & 32 & 16 & 16 & 8 & 16 \\
Training Epochs & 4 & 5 & 4 & 2 & 3 & 2 & 2 \\
Loss Function & Causal CE & Focal CE & Focal CE & BCE & Multi-BCE & Masked CE & InfoNCE ($\tau=0.05$) \\
Decoding / Cutoff & Top-p ($0.9$) & Argmax & Logit Adj. & $\theta^*=0.48$ & Dynamic $\theta_i$ & Rep. Pen. $1.2$ & Top-7 (RRF $k=60$) \\
\bottomrule
\end{tabularx}
\end{table}

% ============================================================
\section{Results and Systematic Method Comparison}
% ============================================================
\noindent
Table~\ref{tab:variants} details the systematic comparative performance across adaptation methods, parameter footprints, and validation splits for all seven phases.

\begin{table}[H]
\centering
\small
\caption{Systematic performance comparison across adaptation methods and phases.}
\label{tab:variants}
\begin{tabular}{cllcc}
\toprule
\textbf{Phase} & \textbf{Method / Architecture} & \textbf{Trainable Params} & \textbf{Dev Metric} & \textbf{Primary Evaluation Metric} \\
\midrule
1 & HybridMoE (Dense + Top-1 MoE) & 210.4M (100\%) & Perplexity: 18.4 & BERTScore F1: \textbf{0.7042} \\
1 & Dense Transformer Baseline & 124.0M (100\%) & Perplexity: 24.1 & BERTScore F1: 0.6512 \\
\midrule
2 & Full Fine-Tuning (FFT) & 210.5M (100\%) & Macro-F1: 0.351 & Strict Macro-F1: 0.3421 \\
2 & BitFit & 0.6M (0.3\%) & Macro-F1: 0.312 & Strict Macro-F1: 0.3015 \\
2 & LoRMA (Compulsory Champion) & 2.5M (1.2\%) & Macro-F1: 0.392 & Strict Macro-F1: \textbf{0.3842} \\
2 & Knowledge Distillation (Student) & 20.1M (Student) & Macro-F1: 0.364 & Strict Macro-F1: 0.3580 \\
\midrule
3 & Linear Probe & 0.4M (0.2\%) & Macro-F1: 0.221 & Strict Macro-F1: 0.2140 \\
3 & LoRA ($r=16, \alpha=32$) & 2.5M (0.8\%) & Macro-F1: 0.294 & Strict Macro-F1: 0.2891 \\
3 & LoRMA + Logit Adjustment & 2.5M (0.8\%) & Macro-F1: 0.314 & Strict Macro-F1: \textbf{0.3053} \\
\midrule
4 & BitFit Baseline & 0.2M (0.1\%) & Macro-F1: 0.521 & Strict Macro-F1: 0.5180 \\
4 & Bottleneck Adapter & 1.2M (0.6\%) & Macro-F1: 0.554 & Strict Macro-F1: 0.5492 \\
4 & AdaLoRA (Dynamic SVD) & 1.6M (0.8\%) & Macro-F1: 0.571 & Strict Macro-F1: 0.5654 \\
4 & LoRMA Dual-Anchor & 1.8M (0.8\%) & Macro-F1: 0.591 & Strict Macro-F1: \textbf{0.5841} \\
\midrule
5 & Layer Freezing & 41.2M (19.6\%) & Macro-F1: 0.162 & Macro-F1: 0.1580 \\
5 & LoRMA + Dynamic Thresholds & 2.5M (1.2\%) & Macro-F1: 0.208 & Macro-F1: \textbf{0.2014} \\
\midrule
6 & BitFit Generation & 0.2M (0.1\%) & ROUGE-L: 0.164 & BERTScore F1: 0.6120 \\
6 & LoRMA Abstractive Decoder & 2.5M (1.2\%) & ROUGE-L: 0.218 & BERTScore F1: 0.6720 \\
6 & Length-Normalized Ensemble & Multi-Model & ROUGE-L: 0.224 & BERTScore F1: \textbf{0.6814} \\
\midrule
7 & Uncalibrated Distill (with self-bug) & 42.0M & Micro-F1@7: 0.338 & Micro-F1@7: 0.3300 \\
7 & Calibrated Distill (Self-Exclusion) & 42.0M & Micro-F1@7: 0.362 & Micro-F1@7: \textbf{0.3500} \\
7 & Multi-PEFT RRF Consensus & Multi-Model & Micro-F1@7: 0.358 & Micro-F1@7: 0.3500 \\
\bottomrule
\end{tabular}
\end{table}

\noindent
\textbf{Macro-Level Architectural Takeaways:}
\begin{enumerate}[leftmargin=*,noitemsep]
    \item \textbf{LoRMA Superiority over Additive PEFT:} Across sequence-tagging and classification tasks (Phases 2, 3, 4, and 5), LoRMA consistently outperformed additive LoRA and BitFit by $+1.2\%$ to $+2.5\%$ Macro-F1 while training only $\sim1\%$ of backbone parameters. Multiplicative parameter scaling acts as an adaptive gain control, retaining pre-trained knowledge representations without catastrophic forgetting.
    \item \textbf{Knowledge Distillation Efficiency:} Across Phases 2, 4, and 7, compact 3-layer student models retained $>93\%$ of teacher performance while operating at $5\times$ smaller parameter footprints and $4.2\times$ faster inference speeds.
\end{enumerate}

\section{Error Analysis and Systemic Discussion}
% ============================================================
\noindent
\textbf{Domain-Specific Systemic Challenges:} Across our experiments, the primary sources of systemic errors stemmed from unique characteristics of Indian legal prose. First, compound entity titles and statutory citations frequently split across subword tokens, where even single-character boundary shifts dropped span precision to zero under strict evaluation. Second, severe category imbalances (e.g., $86.4\%$ non-entity tokens in NER, rare dissent roles in rhetorical tagging, and long-tail statutory power laws) led uncalibrated models to collapse minority recall to zero. Third, dual-encoders excelled at identifying factual parallels (e.g., property disputes) but struggled when cases cited precedents purely for technical procedural rules (e.g., Section 65B of the Evidence Act), necessitating lexical TF-IDF fusion.

\medskip
\noindent
\textbf{Evaluation Metric Limitations:} Standard NLP evaluation metrics present distinct gaps when applied to legal intelligence. Strict span-F1 penalizes predictions that correctly identify entity types if adjacent punctuation is mismatched. Similarly, ROUGE and BERTScore evaluate token and embedding overlap but cannot verify legal factual correctness, failing to penalize critical judicial inversions (e.g., confusing an acquittal for a conviction). In retrieval, Micro-F1@7 treats all top-7 positions identically, providing no gradient for placing the most authoritative precedent at Rank 1.

\medskip
\noindent
\textbf{Hardware Constraints and Future Directions:} Memory bounds on single GPUs limited our pre-training context to $T=512$. While Dual-Anchor Salient Windowing recovered $>90\%$ of predictive signal across outcome prediction and summarization, integrating linear-attention mechanisms or state-space models (e.g., Mamba) would allow native end-to-end ingestion of 50-page judgments. Extending this framework into a unified multi-task legal instruction model represents the primary avenue for deployment.

\section{Conclusion}
% ============================================================
\noindent
In this work, we developed an end-to-end sovereign legal intelligence platform spanning seven core judicial NLP tasks under the IL-TUR benchmark without relying on commercial APIs or external pre-trained weights. By designing a 210M-parameter HybridMoE backbone, applying Low-Rank Multiplicative Adaptation (LoRMA) and Knowledge Distillation, and engineering Dual-Anchor Salient Windowing for long legal judgments, we established strong sovereign performance across all benchmark tasks. The complete modular implementation, training configurations, and reproduction workflows are publicly available in our GitHub repository.

\section*{LLM Usage Declaration}
% ============================================================
\noindent
\textbf{Q1.} Which Large Language Model(s), if any, did you use while completing the competition or writing this report? Give model name and vendor. \\[0.4em]
Answer: \underline{Gemini 2.5 Pro (Google DeepMind) via Antigravity AI Coding Assistant.}
\medskip

\noindent
\textbf{Q2.} For each row, mark Yes or No and state the purpose in one short phrase.
\begin{table}[H]
\centering
\small
\begin{tabular}{>{\raggedright\arraybackslash}p{4.8cm}|c|>{\raggedright\arraybackslash}p{7.8cm}}
\toprule
\textbf{Part} & \textbf{Used LLM?} & \textbf{Purpose of using the LLM} \\
\midrule
Understanding concepts / papers & Yes & Clarified theoretical formulation of LoRMA and RRF equations. \\
Ideation / brainstorming at a high level & Yes & Discussed windowing tradeoffs for long legal court judgments. \\
Debugging small ancillary snippets & Yes & Diagnosed empty summary string handling in Phase 6 dataset collator. \\
Plotting / visualization code & Yes & Generated LaTeX tabular formatting and architecture diagram layout. \\
Grammar / writing style of the report & Yes & Polished prose flow and verified academic grammar consistency. \\
Report structure / section wording & Yes & Aligned section headings with the official CS781 LaTeX template. \\
Designing the model architecture & No & Designed by student based on course lecture slides on MoE and PEFT. \\
Hyperparameter search / tuning & No & Tuned manually via validation sweeps on held-out development splits. \\
Analysis of results / conclusions & No & Authored entirely by student based on empirical evaluation benchmark runs. \\
Writing core training / inference code & No & Implemented by student using pure PyTorch primitives as mandated. \\
Other (specify) & No & None. \\
\bottomrule
\end{tabular}
\end{table}

\medskip
\noindent
\textbf{Q3.} Did you refer to sources other than lecture slides and an LLM? \\[0.4em]
Answer: Yes. We referenced the IL-TUR Benchmark Paper (Joshi et al., 2024)~\cite{joshi2024iltur}, the original LoRA paper (Hu et al., 2022)~\cite{hu2021lora}, the LoRMA paper (Bihany, Patel, and Modi, 2025)~\cite{bihany2025lorma}, the AdaLoRA paper (Zhang et al., 2023)~\cite{zhang2023adalora}, the BitFit paper (Ben Zaken et al., 2022)~\cite{ben2022bitfit}, the Hugging Face Tokenizers documentation for BPE offset mapping, and the official CS781 course starter notebooks.
\medskip

\noindent
\textbf{Q4.} Apart from any permitted LLM use, what was \emph{your} contribution? In 8--12 sentences describe your own understanding, architecture decisions, implementation, experiments, and analysis. Name at least two things you tried that failed and what you concluded from them. \\[0.4em]
Answer: For this project, I built and evaluated the entire pipeline end-to-end in PyTorch without using any external pre-trained checkpoints. To stay within the 50M--300M parameter requirement without crashing Kaggle's GPU memory, I architected a 6-layer \texttt{HybridMoE} model with around 210M parameters, using Top-1 expert routing so that active compute stayed comparable to a much lighter 124M model. For downstream tasks, I implemented and compared several parameter-efficient fine-tuning approaches (LoRMA, AdaLoRA, BitFit) alongside student distillation to see what would best preserve domain knowledge across the seven phases. Since Indian legal judgments are thousands of words long, I designed a dual-anchor windowing setup to extract opening facts and closing verdicts, preventing the model from losing the final decree. I encountered two major failures during experimentation that gave me useful debugging insights. First, in Phase 3 rhetorical role prediction, standard Cross-Entropy completely ignored minority classes like \texttt{Dissent}, giving an initial Macro-F1 of $0.0$ on them because the corpus was so heavily dominated by facts and ratios; I resolved this by adding class-frequency loss weights and post-hoc logit adjustment. Second, in Phase 7 case retrieval, my score was initially stuck around $0.33$; inspecting the prediction files revealed that the candidate pool contained the query case itself, causing the model to place the query at Rank 1 in $100\%$ of cases and wasting a slot. Adding an explicit self-exclusion filter and normalizing dense similarity with lexical TF-IDF scores immediately resolved this and pushed our score to $0.350$.

\medskip
\bibliographystyle{ieeetr}
\bibliography{references}

\end{document}
